\chapter{Đọc thêm 3: Phân tích dữ liệu vòng đệm O-ring của Challenger (Wujek)}
\label{ch:M6L4DocThem3}

\section{Thông tin nguồn}

\begin{itemize}
  \item \textbf{Nguồn:} Wujek, J. H. ``Challenger O-Ring Data Analysis.'' Numerical \& Design Problems With Ethical Content, TAMU, 1995.
  \item \textbf{Phần cần đọc:} Toàn bộ
  \item \textbf{Ghi chú:} bản chuyển ngữ tiếng Việt súc tích, bám cấu trúc nguồn (giữ nguyên tiêu đề, công thức, số liệu, bảng, chú thích và danh mục tài liệu tham khảo).
\end{itemize}

\textbf{Tác giả:} Joseph H. Wujek. \textbf{Bộ sưu tập:} Numerical \& Design Problems With Ethical Content. \textbf{Đơn vị:} Zachry Department of Civil Engineering-TAMU Ethics. \textbf{Năm:} 1995. \textbf{DOI:} https://doi.org/10.18130/2pem-7a20

\textbf{Lĩnh vực:} Kỹ thuật hàng không vũ trụ; Khoa học máy tính, toán và vật lý; Kỹ thuật; Khoa học và kỹ thuật vật liệu; Kỹ thuật cơ khí; Thống kê và xác suất.

\textbf{Chủ đề:} Thảm họa, hiểm họa; Quan hệ giữa người sử dụng lao động và người lao động; An toàn phòng thí nghiệm và nơi làm việc; An toàn; Đạo đức nghề nghiệp. Loại: nghiên cứu tình huống.

\subsection{Mô tả}

Tình huống này phân tích dữ liệu vòng đệm O-ring trong thảm họa Challenger và lập luận cho việc hủy bỏ lần phóng. Phù hợp với các môn thống kê ứng dụng, vật liệu và kỹ thuật đại cương (cấp 1-4).

\subsection{Giới thiệu}

Các sự kiện dẫn đến thảm họa Challenger đã được biết rõ. Bài toán này phân tích dữ liệu vòng đệm O-ring và dựa trên kết quả để lập luận về việc hủy phóng.

Dữ liệu ở Bảng 1 được David Hodges sử dụng trong Hội thảo Sinh viên năm nhất tại Đại học California, Berkeley.¹ Đây là dữ liệu của các chuyến bay thành công. Roger Boisjoly cho biết các dữ liệu này chưa có trước thảm họa Challenger.² Hãy dùng dữ liệu Bảng 1 khi cần.³

\begin{enumerate}
\def\labelenumi{\arabic{enumi}.}
\item
  Lập đồ thị hồi quy tuyến tính và ngoại suy kết quả để dự đoán mức hao mòn vật liệu của vòng đệm O-ring tại thời điểm phóng Challenger thảm họa. Giả sử vòng đệm ở nhiệt độ môi trường bệ phóng, 29 °F.

  \begin{enumerate}
  \def\labelenumii{\alph{enumii}.}
    \item
    Tính hệ số tương quan r. So sánh với giá trị tới hạn của r ở mức ý nghĩa 5\% và giải thích ý nghĩa.
  \item
    Bình luận về ý nghĩa của r, \(r^2\) và \(1-r^2\). Điều này có liên quan thế nào khi lập luận không nên phóng?
  \item
    Tính cận khoảng tin cậy 95\% và thể hiện kết quả trên đồ thị tuyến tính.
  \item
    Bình luận về độ phù hợp bằng kiểm định Chi-square (khi bình phương).
  \item
    Bình luận về độ phù hợp bằng kiểm định Kolmogorov-Smirnov.
  \end{enumerate}
\end{enumerate}

\textbf{Bảng 1. Dữ liệu đo sau chuyến bay của Challenger³}

{\small
\begin{longtable}[]{@{}ll@{}}
\toprule\noalign{}
Nhiệt độ O-ring (°F) & Độ sâu xói mòn \ensuremath{\delta} (mil) * \\
\midrule\noalign{}
\endhead
\bottomrule\noalign{}
\endlastfoot
66.0 & 0.0 \\
70.0 & 53.0 \\
69.0 & 0.0 \\
68.0 & 0.0 \\
67.0 & 0.0 \\
72.0 & 0.0 \\
73.0 & 0.0 \\
70.0 & 0.0 \\
57.0 & 40.0 \\
63.0 & 0.0 \\
70.0 & 28.0 \\
78.0 & 0.0 \\
67.0 & 0.0 \\
53.0 & 48.0 \\
67.0 & 0.0 \\
75.0 & 0.0 \\
70.0 & 0.0 \\
81.0 & 0.0 \\
76.0 & 0.0 \\
79.0 & 0.0 \\
75.0 & 0.0 \\
76.0 & 0.0 \\
\end{longtable}
}

* 1 mil bằng \(10^{-3}\) inch.

\begin{enumerate}
\def\labelenumi{\arabic{enumi}.}
\setcounter{enumi}{1}
\item
  Thử khớp các đa thức bậc n = 2, 3, ... với dữ liệu và áp dụng các thước đo ở bài 1(c), (d) và (e) cho kết quả.
\item
  Thử biến đổi tập dữ liệu bằng hàm logarit hoặc hàm lũy thừa rồi khớp hàm tuyến tính với các biến đã biến đổi. Áp dụng các thước đo ở bài 1(c), (d) và (e) cho kết quả và so sánh.
\item
  Với các kết quả trên, dữ liệu sẽ hữu ích đến mức nào khi lập luận cho quyết định "không phóng"? Nhận xét về các hệ quả đạo đức.
\end{enumerate}

\subsection{LỜI GIẢI}

Giả định sinh viên đã được học về các khía cạnh kỹ thuật của thảm họa.

Mục đích của các bài tập là giúp sinh viên nhạy bén với các vấn đề sau:

\begin{itemize}
\item
  Thế giới mang tính xác suất, không tất định, và không có gì là "chắc chắn".
\item
  Dù dữ liệu có thể chưa đủ kết luận, khi tính mạng và tài sản bị đe dọa thì "không biết" phải được coi là "chúng ta đang có vấn đề".
\item
  Hồi quy tuyến tính chỉ là một cách khớp hàm với dữ liệu và diễn giải kết quả về mặt thống kê.
\end{itemize}

\begin{enumerate}
\def\labelenumi{\arabic{enumi}.}
\item
  (Đồ thị được cho bên dưới, phù hợp để làm slide chiếu.) (Phương trình ghi ba chữ số có nghĩa để đối chiếu kết quả.) Tại 29 °F, mức hao hụt vật liệu (độ tin cậy 95 \%) nằm trong khoảng 25 đến 105 mil. Biến thiên cực lớn này cho thấy (xem thêm bên dưới) rủi ro hỏng gioăng rất cao. Dù ngoại suy luôn đáng ngờ, chính việc "không biết" một cách có ý nghĩa thống kê gioăng sẽ hoạt động ra sao ở 29 °F lẽ ra đã là lý do đủ để hủy phóng.

  \begin{enumerate}
  \def\labelenumii{\alph{enumii}.}
  \item
    Từ Hình 1, r = \ensuremath{-}0,56. Tra bảng giá trị tới hạn của hệ số tương quan, ở mức tin cậy 95 \% (mức ý nghĩa 5 \%), với 20 bậc tự do {[}22 cặp dữ liệu trừ hai tham số của phương trình tuyến tính{]}, tìm được \(r_c = 0{,}423\). Điều này có nghĩa là với xác suất 5 \%, dữ liệu có hệ số tương quan lớn đến 0,423 vẫn có thể là không tương quan. Hay nói cách khác, \(r_c\) là giá trị lớn nhất mà \(r\) có thể đạt được chỉ do ngẫu nhiên (trong 95 \% số lần) khi không tồn tại tương quan.
  \item
    Hệ số tương quan r được định nghĩa là: \(r = S_{XY}/(S_X S_Y)\). Tử số là hiệp phương sai mẫu, mẫu số là tích của các độ lệch chuẩn mẫu của hai biến X, Y. Do đó giá trị r là thước đo mức độ liên hệ giữa hai biến. Ta có \(0 \le |r| \le 1\), với \(r = \pm 1\) là tương quan hoàn hảo và \(r = 0\) là hoàn toàn không tương quan. Khi đó \(r^2\) là thước đo phần biến thiên do quan hệ nhân quả, còn \(1 - r^2\) là thước đo phần biến thiên do ngẫu nhiên. Xem thêm các nhận xét ngay sau đề bài ở trên.
  \item
    Xem Hình 1. Giải thích ý nghĩa của khoảng tin cậy. Một cách diễn giải: "Nếu thực hiện một số rất lớn phép thử, 95 \% kết quả sẽ nằm trong khoảng giới hạn 95 \% đã chỉ ra."
  \item
    Đối với cả kiểm định Chi-square và kiểm định K-S ở bài (e), giả thuyết kiểm định là: \(H_0\): dữ liệu không tương quan.

    Ta chấp nhận giả thuyết nếu: \(\chi_0^2 > \chi^2_{1-\alpha,\nu}\)

    trong đó \(1 - \alpha = C\) là mức tin cậy (\(\alpha\) là mức ý nghĩa), và \(\nu\) là số bậc tự do. Ở đây, \(\nu = N - 1 - m\), với N là số quan sát của các biến (22 trong trường hợp này) và m là số tham số của phương trình được kiểm định (m = 2 với phương trình tuyến tính).

    Vế trái của bất đẳng thức là thống kê Chi-square quan sát được, tính từ dữ liệu. Vế phải là giá trị của biến Chi-square ứng với xác suất C, tức là giá trị của tích phân hàm mật độ đến biến ngẫu nhiên đó; các giá trị này đã được lập bảng.

    Với bài toán đang xét: \(\chi^2_{0.95,19} = 30{,}1\) và \(\chi_0^2 = 422\).

    Do đó, giả thuyết được chấp nhận: dữ liệu không tương quan tuyến tính.
  \item
    Với kiểm định Kolmogorov-Smirnov, ta dùng cùng giả thuyết: \(H_0\): dữ liệu không tương quan.

    Ta bác bỏ giả thuyết nếu: \(\sup\{|y_f - y_d|\} \ge d_\alpha(\nu)\)

    Chỉ số dưới của y lần lượt là giá trị "khớp" (fitted) và giá trị "dữ liệu" (data). Như vậy, với một cặp dữ liệu \(x_d, y_d\), giá trị của hàm khớp tại \(x_d\) là \(y_f\). Giá trị lớn nhất ("supremum") được so sánh với thống kê K-S ứng với mức ý nghĩa và số bậc tự do như ở bài (d).

    Ta có \(\sup\{|y_f - y_d|\} = 45\) và \(d_{0.05}(19) = 0{,}301\), nên không thể bác bỏ giả thuyết rằng dữ liệu không tương quan tuyến tính.
  \end{enumerate}
\item
  Bài toán 2 dành cho các bài tập khớp đường cong. Do tính chất mở của bài, không cung cấp lời giải. Một cạm bẫy khi dùng phép biến đổi logarit là cố tính logarit của số không!
\item
  Bài toán 3 dành cho các bài tập khớp đường cong. Do tính chất mở của bài, không cung cấp lời giải. Một cạm bẫy khi dùng phép biến đổi logarit là cố tính logarit của số không!
\item
  Dữ liệu thất thường lẽ ra đã đủ để báo hiệu "chúng ta đang có vấn đề" vì rủi ro đối với tính mạng con người. Các phi hành gia không được thông báo về lịch sử các khiếm khuyết của vòng đệm O-ring, nên không thể có sự đồng thuận có hiểu biết.
\end{enumerate}

\subsection{Chú thích}

\begin{enumerate}
\def\labelenumi{\arabic{enumi}.}
\item
  Lighthall, F. "Launching the Space Shuttle Challenger: Disciplinary Deficiencies in the Analysis of Engineering Data," IEEE Transactions on Engineering Management, Vol. 38, No. 1, February 1991. pp. 63-74.
\item
  Boisjoly, R. Pers. comm. November 1993.
\item
  Lighthall, p. 70, Table I.
\item
  Found in many texts and handbooks. See for example, Fisher, R., and Yates, F., Statistical Tables for Biological. Agricultural, and Medical Research. Oliver \& Boyd, Ltd., Edinburgh, 1957.
\item
  Owen, D. Handbook of Statistical Tables. Addison-Wesley, Reading, MA, 1962. p. 51.
\item
  Owen, p. 64.
\end{enumerate}

Hình 1. Biểu đồ hồi quy tuyến tính trên dữ liệu vòng đệm O-ring (theo Lighthall)

\subsection{Ghi chú}

Tác giả: Dr. Joseph H. Wujek, P.E., College of Engineering, University of California, Berkeley.

Các bài toán này ban đầu được xây dựng trong khuôn khổ một dự án do NSF tài trợ nhằm tạo ra các bài toán số có chứa vấn đề đạo đức, dùng cho bài tập của các môn kỹ thuật và môn học khác. Các bài toán trình bày ở đây đã được chỉnh sửa đôi chút cho rõ ràng.

\subsection{Trích dẫn}

Joseph H. Wujek. 1995. Challenger O-Ring Data Analysis. Online Ethics Center. DOI:https://doi.org/10.18130/2pem-7a20.
